\originalchapter{8.05 作业2：节点、变分与四次势}
来源为Fall 2013 Assignment 2, 2013年9月13日发布, 9月20日15时截止. 其中正文页眉沿用Fall 2008, 不改网页托管学期. 第1--6题中文转述与\textbf{AI独立解答}; 第4题保留Shankar改编署名, 使用MIT完整给出的任务独立证明.
\originalsection{第1题：中央点障壁的近简并}
\begin{exercise}
无限井$0<x<a$中央加$V_0a\delta(x-a/2)$, $\gamma=mV_0a^2/\hbar^2\gg1$. 求基态能量至$1/\gamma$阶, 与第一激发态比较并解释间隙.
\end{exercise}
\begin{solution}
关于$a/2$奇的态在障壁处为0, 最低奇态是原$n=2$驻波, 能量$E_o=2\pi^2\hbar^2/(ma^2)$. 最低偶态左半井$C\sin kx$, 右边对称, 跳跃条件给$-2k\cot(ka/2)=2mV_0a/\hbar^2$. 令$z=ka/2=\pi-\varepsilon$, 得$z\cot z=-\gamma/2$, $\varepsilon=2\pi/\gamma+O(\gamma^{-2})$. 所以$E_g=E_o(1-4/\gamma+O(\gamma^{-2}))$, 间隙$E_o-E_g=8\pi^2\hbar^2/(ma^2\gamma)+O(\gamma^{-2})$. 无限高障壁将左右两长度$a/2$的井隔开, 对称和反对称组合趋简并; 有限障壁仍有隧穿, 不违反一维束缚谱非简并.
\end{solution}
\originalsection{第2题：相邻节点之间}
\begin{exercise}
$\psi_j''+(E_j-U)\psi_j=0$, $E_{k+1}>E_k$. 推$[\psi_{k+1}\psi_k'-\psi_k\psi_{k+1}']_a^b=(E_{k+1}-E_k)\int_a^b\psi_k\psi_{k+1}\dd x$. 若$a,b$为$\psi_k$相邻零点且其中正, 证$\psi_{k+1}$在其中变号.
\end{exercise}
\begin{solution}
左侧括号求导, 混合一阶项抵消, 利用两方程得$(E_{k+1}-E_k)\psi_k\psi_{k+1}$, 积分即成. 假定高态不变号并取其非负, 正规非零解若有内零且不变号必函数与导数同为0, 与唯一性矛盾, 所以内部严格正. 低态有$\psi_k'(a)>0$, $\psi_k'(b)<0$, 左侧$\psi_{k+1}(b)\psi_k'(b)-\psi_{k+1}(a)\psi_k'(a)\le0$, 右侧严格正, 矛盾. 因而至少一内节点. 这里假定普通Sturm--Liouville问题的连续势和正规边界, 不把端点无穷或奇异节点一概视同有限区间.
\end{solution}
\originalsection{第3题：变分的鞍点}
\begin{exercise}
(a)归一试态与基态正交, 证能量均值不低于第一激发能. (b)实本征基中$\psi=\psi_2+\sum_n\varepsilon_n\psi_n$, 展Rayleigh商至二阶, 哪些方向消失? 是极大、极小还是鞍点?
\end{exercise}
\begin{solution}
展开$\psi=\sum_{n\ge2}c_n\psi_n$, $\sum|c_n|^2=1$, 平均$\sum E_n|c_n|^2\ge E_2$. 这要求正交约束准确, 数值污染基态会破坏界. 未约束商的精确式为
\[
F=\frac{E_2(1+\varepsilon_2)^2+\sum_{n\ne2}E_n\varepsilon_n^2}
{(1+\varepsilon_2)^2+\sum_{n\ne2}\varepsilon_n^2}
=E_2+\sum_{n\ne2}(E_n-E_2)\varepsilon_n^2+O(\varepsilon^3).
\]
所有一次项抵消, $\varepsilon_2$是仅变整体尺度的平坦方向. $n=1$负曲率, 高能方向正曲率, 故在至少3能级时为鞍点; 若限定基态正交子空间, 则成为最小. 有简并再出现对应平坦方向.
\end{solution}
\originalsection{第4题：任意一维吸引势为何能束缚}
\begin{exercise}
$V\le0$, 分段连续并在无穷远趋0, 在某个连续点严格负. 以归一Gaussian$\psi_\alpha=(\alpha/\pi)^{1/4}\e^{-\alpha x^2/2}$证明至少有一个负能束缚态. 原件署based on Shankar.
\end{exercise}
\begin{solution}
连续性给有限区间$[x_1,x_2]$上$V\le-v_0<0$. 令$d=x_2-x_1$, $R=\max(|x_1|,|x_2|)$. Gaussian动能为$\hbar^2\alpha/(4m)$, 其余区间势能非正, 因而
\[
\langle H\rangle\le\frac{\hbar^2\alpha}{4m}
-v_0d\sqrt{\frac{\alpha}{\pi}}\e^{-\alpha R^2}<0
\quad(\alpha\text{充分小}).
\]
负项为$\sqrt\alpha$阶, 比正动能的$\alpha$阶降得慢. 因此宽试态虽在井内只留很少概率, 仍能降低总能量. 对有下界且趋0的普通短程势, 本质谱下沿为0; 变分原理给其下方一个离散本征值. 这一步是由负Rayleigh商到束缚态的谱依据. 原题的“非零”须在几乎处处意义下理解或具有上述负值连续点; 仅改变一个孤立点的势值不会改变Hamilton算符, 不能推出束缚.
\end{solution}
\originalsection{第5题：四次势的数值与变分}
\begin{exercise}
$H=-\hbar^2\partial_x^2/(2m)+\alpha x^4$, $\alpha>0$. (a)令$x=\beta u$化为$-\psi''/2+u^4\psi=e\psi$. (b)偶初值$\psi(0)=1,\psi'(0)=0$求基态$e$至六位. (c)构造第一激发能的变分上界. (d)数值第一激发能与界比较.
\end{exercise}
\begin{solution}
动能与势系数要求$\alpha\beta^4=\hbar^2/(m\beta^2)$, 所以$\beta=(\hbar^2/(m\alpha))^{1/6}$, 能量尺度$\gamma=(\hbar^4\alpha/m^2)^{1/3}$, $E=\gamma e$. 偶基态在原点导数0. 采用Hermite振子基的Galerkin矩阵, 增大截断维数并核对低能谱收敛, 得$e_0=0.66798626$, 第一奇态$e_1=2.39364402$（六位分别0.667986、2.393644）.

为保证与偶基态正交, 取奇试态$\psi_b=Nu\e^{-bu^2/2}$,$b>0$. Gaussian矩给$\langle T\rangle=3b/4$, $\langle u^4\rangle=15/(4b^2)$, 所以$e_1\le3b/4+15/(4b^2)$. 驻点$b^3=10$为全局极小, 上界$9\sqrt[3]{10}/8\simeq2.42373903$. 数值能量低于此界约1.26\%. 试函数奇性至关重要, 任意Gaussian只能保证基态上界. 数值逼近是谱计算证据, 不替代上界的数学证明; 自有Python脚本实现, 未复制商业软件教材代码.
\end{solution}
\originalsection{第6题：矩阵指数}
\begin{exercise}
$M=\begin{pmatrix}0&-\ii\\\ii&0\end{pmatrix}$, 求$\e^{\ii M\theta}=A(\theta)I+B(\theta)M$. 哪种任意阶矩阵代数性质可给同结果?
\end{exercise}
\begin{solution}
$M^2=I$, 所以偶次幂均$I$, 奇次幂均$M$. Taylor级数分组为$I\cos\theta+\ii M\sin\theta$, $A=\cos\theta$, $B=\ii\sin\theta$. 任意尺寸满足$M^2=I$便有同式, 不要求只有两个维度. 若同时Hermitian则指数幺正; 仅平方为$I$能给代数公式, 不自动给内积保持.
\end{solution}
